\begin{appendix}

\section{Dirichlet aggregation and the moments of linear functionals}\label{app:dirichlet}

Let $\bpi\sim\Dir(\balpha)$ on the $J$-simplex, with $\alpha_0=\sum_j\alpha_j$, and for $\cR\subseteq\{1,\dots,J\}$ write $\pi(\cR)=\sum_{j\in\cR}\pi_j$ and $\alpha(\cR)=\sum_{j\in\cR}\alpha_j$.

\begin{lemma}[Aggregation]\label{lem:agg}
For any partition $\cR_1,\dots,\cR_L$ of $\{1,\dots,J\}$, $\big(\pi(\cR_1),\dots,\pi(\cR_L)\big)\sim\Dir\big(\alpha(\cR_1),\dots,\alpha(\cR_L)\big)$. In particular $\pi(\cR)\sim\Beta\big(\alpha(\cR),\alpha_0-\alpha(\cR)\big)$ for every $\cR$, which is \eqref{eq:beta}.
\end{lemma}

\begin{proof}
Write $\pi_j=G_j/\sum_{l}G_l$ with $G_1,\dots,G_J$ independent, $G_j\sim\mathrm{Gamma}(\alpha_j,1)$. Then $\sum_{j\in\cR_l}G_j\sim\mathrm{Gamma}(\alpha(\cR_l),1)$, the $L$ group sums are independent, and the vector of their proportions is Dirichlet with the group-sum parameters. The Beta statement is the case $L=2$.
\end{proof}

\begin{lemma}[Moments of linear functionals]\label{lem:mom}
For $\bc\in\mathbb R^J$, $\E[\bc^\top\bpi]=\bc^\top\balpha/\alpha_0$ and
\begin{equation*}
\Var[\bc^\top\bpi]=\frac{\alpha_0\sum_j c_j^2\alpha_j-\big(\sum_j c_j\alpha_j\big)^2}{\alpha_0^2(\alpha_0+1)} .
\end{equation*}
\end{lemma}

\begin{proof}
The Dirichlet moments are $\E[\pi_j]=\alpha_j/\alpha_0$ and $\Cov(\pi_j,\pi_l)=(\alpha_0\,\alpha_j\ind\{j=l\}-\alpha_j\alpha_l)/\{\alpha_0^2(\alpha_0+1)\}$, and $\Var[\bc^\top\bpi]=\sum_{j,l}c_jc_l\Cov(\pi_j,\pi_l)$.
\end{proof}

Equation \eqref{eq:moments} follows with $\bc(\bw)=\sum_k w_k\ind_{\cR_k}$ applied to the posterior $\Dir(\balpha'_a)$: $\sum_j c_j\alpha'_{a,j}=\sum_k w_k\alpha'_a(\cR_k)$, and since $\ind_{\cR_k}(j)\ind_{\cR_{k'}}(j)=\ind_{\cR_k\cap\cR_{k'}}(j)$, $\sum_j c_j^2\alpha'_{a,j}=\sum_{k,k'}w_kw_{k'}\alpha'_a(\cR_k\cap\cR_{k'})$. The same calculation with $\bc=\ind_{\cR_k}$ and $\bc=\ind_{\cR_{k'}}$ gives the posterior covariance of two component success probabilities, $\Cov(p_{a,k},p_{a,k'}\mid\bx^{(1)}_a)=\{\alpha'_{a,0}\alpha'_a(\cR_k\cap\cR_{k'})-\alpha'_a(\cR_k)\alpha'_a(\cR_{k'})\}/\{(\alpha'_{a,0})^2(\alpha'_{a,0}+1)\}$, which is the quantity the conjunctive gate and the composite gate share.

\section{Proof of Theorem~\protect\thmtone}\label{app:proof}

Throughout, $\mathcal F_1$ denotes the $\sigma$-field generated by the stage-1 data and the external information used at the interim; the selected dose $a^\star$, the go/no-go decision, $n_2$ and every $p^{(1)}_{a,k}$ are $\mathcal F_1$-measurable. A statistic $P$ with values in $[0,1]$ is called a valid $p$-value for a hypothesis $H$ if $\Pr(P\le u)\le u$ for all $u\in[0,1]$ whenever $H$ is true, and conditionally valid given $\mathcal F_1$ if $\Pr(P\le u\mid\mathcal F_1)\le u$ almost surely on the event that $H$ is true.

\begin{lemma}[Combination of valid $p$-values]\label{lem:comb}
Let $\psi_m$ be the Bonferroni combination, or the Simes combination applied to $p$-values satisfying the Simes inequality. If $P_1,\dots,P_m$ are valid $p$-values for $H_1,\dots,H_m$, then $\psi_m(P_1,\dots,P_m)$ is a valid $p$-value for $\bigcap_iH_i$. The same holds conditionally on $\mathcal F_1$: if the $P_i$ are conditionally valid and satisfy the Simes inequality conditionally, $\psi_m(P_1,\dots,P_m)$ is conditionally valid. Both combinations are non-decreasing in each argument.
\end{lemma}

\begin{proof}
For Bonferroni, $\Pr(m\min_iP_i\le u)\le\sum_i\Pr(P_i\le u/m)\le u$, and the same inequality holds for conditional probabilities. For Simes the statement is the Simes inequality\cite{simes1986}, proved by Sarkar\cite{sarkar1998} for multivariate totally positive statistics and by Benjamini and Yekutieli\cite{benjamini2001} under positive regression dependence on the subset of true hypotheses, which includes jointly normal statistics with non-negative correlations; applied to the conditional distribution given $\mathcal F_1$ it gives the conditional statement. Monotonicity is immediate from the formulas, since increasing any argument cannot decrease any order statistic.
\end{proof}

\begin{lemma}[Intersection--union $p$-value]\label{lem:iu}
If $p_1,\dots,p_K$ are valid $p$-values for $H_1,\dots,H_K$, then $q=\max_k p_k$ is a valid $p$-value for $\bigcup_k H_k$.
\end{lemma}

\begin{proof}
If $\bigcup_kH_k$ is true, some $H_{k_0}$ is true, and $\Pr(q\le u)\le\Pr(p_{k_0}\le u)\le u$.
\end{proof}

\begin{lemma}[Inverse-normal combination across stages]\label{lem:inv}
Let $P_1$ be $\mathcal F_1$-measurable and valid, and let $P_2$ be conditionally valid given $\mathcal F_1$. Then for $v_1^2+v_2^2=1$, $v_2>0$, the combination $\cC(P_1,P_2)$ of \eqref{eq:comb} satisfies $\Pr\big(\cC(P_1,P_2)\le\alpha\big)\le\alpha$.
\end{lemma}

\begin{proof}
Put $Z_i=\Phi^{-1}(1-P_i)$, so that $\cC(P_1,P_2)\le\alpha$ if and only if $v_1Z_1+v_2Z_2\ge z_{1-\alpha}$. Conditional validity gives $\Pr(Z_2\ge t\mid\mathcal F_1)=\Pr\big(P_2\le1-\Phi(t)\mid\mathcal F_1\big)\le1-\Phi(t)$ for all $t$, hence
\begin{equation*}
\Pr\big(\cC\le\alpha\big)=\E\Big[\Pr\Big(Z_2\ge\frac{z_{1-\alpha}-v_1Z_1}{v_2}\,\Big|\,\mathcal F_1\Big)\Big]\le\E\big[g(Z_1)\big],\qquad g(z)=1-\Phi\Big(\frac{z_{1-\alpha}-v_1z}{v_2}\Big),
\end{equation*}
with $g$ non-decreasing. Validity of $P_1$ gives $\Pr(Z_1\ge t)\le1-\Phi(t)$, so $Z_1$ is stochastically smaller than a standard normal variable $Z$, and $\E[g(Z_1)]\le\E[g(Z)]=\Pr(v_1Z+v_2Z'\ge z_{1-\alpha})=\alpha$ for $Z'$ standard normal and independent of $Z$.
\end{proof}

\begin{proof}[Proof of Theorem~\ref{thm:t1e}]
Fix a configuration $(\bpi_1,\dots,\bpi_M,\bpi_C)$ and let $T=\{(a,k)\in\cJ:H_{a,k}\text{ true}\}$ be the set of true elementary hypotheses. It is a fixed set, determined by the configuration alone; if it is empty there is nothing to prove. By the closure principle, the rejection of any true $H_{a,k}$, $(a,k)\in T$, requires the rejection of every intersection containing $(a,k)$, in particular of $H_T$; a false claim rejects some true $H_{a^\star,k_0}$ and is therefore such a rejection. It suffices to show that $\Pr(H_T\text{ rejected})\le\alpha$.

A trial that stops for futility rejects nothing, so the rejection event is contained in the corresponding event for the trial that always proceeds to stage 2; we work with the latter, in which stage-2 data exist on every sample path and (i) and (iii) hold as stated. By \eqref{eq:rejJ}, $H_T$ is rejected if and only if $\cC(P_1,P_2)\le\alpha$ with
\begin{equation*}
P_1:=p^{(1)}_T=\psi_{|T|}\big((p^{(1)}_{a,k})_{(a,k)\in T}\big),\qquad
P_2:=p^{(2)}_T=\psi_{|T^\star|}\big((p^{(2)}_k)_{k\in T^\star}\big),\quad T^\star=\{k:(a^\star,k)\in T\},
\end{equation*}
and $P_2=1$ when $T^\star=\emptyset$. The stage-1 component $P_1$ is $\mathcal F_1$-measurable, and it is a valid $p$-value for $H_T$ by (ii) and Lemma~\ref{lem:comb}, because every $p^{(1)}_{a,k}$ entering it tests a true hypothesis. Given $\mathcal F_1$, the selected dose $a^\star$ and hence the set $T^\star$ are fixed. On $\{T^\star\ne\emptyset\}$ every $H_{a^\star,k}$ with $k\in T^\star$ is true, so by (i) and (iii) the $p^{(2)}_k$, $k\in T^\star$, are conditionally valid given $\mathcal F_1$ and, with the conditional Simes inequality where the Simes combination is used, Lemma~\ref{lem:comb} makes $P_2$ conditionally valid. On $\{T^\star=\emptyset\}$, $P_2=1$ is trivially conditionally valid. Lemma~\ref{lem:inv} gives $\Pr(\cC(P_1,P_2)\le\alpha)\le\alpha$, which completes the proof.
\end{proof}

The argument uses the selection rule, the futility rule and $n_2$ only through their $\mathcal F_1$-measurability. What the closure over all $MK$ pairs buys is that the true intersection $H_T$ is fixed in advance: it contains each dose's own null endpoints, whichever they are, so the stage-1 component $P_1$ is a single valid $p$-value whatever dose is selected. A closure over doses on a fixed endpoint $k$ would use $\bigcap_{a}H_{a,k}$ over the doses null on $k$, a set that changes with the endpoint on which the selected dose happens to be null; Appendix~\ref{app:counter} shows what goes wrong.

\subsection{The dose-level alternative}\label{app:union}

The dose-level closure of Section~\ref{sec:final} tests the hypotheses $H^\cup_a$, $a=1,\dots,M$, by closed testing over doses, each intersection $H^\cup_I=\bigcap_{a\in I}H^\cup_a$, $I\ni a^\star$, being rejected if and only if $\cC\big(\psi_{|I|}(p^{(1)}_{a^\star,k},(q^{(1)}_a)_{a\in I\setminus\{a^\star\}}),\,p^{(2)}_k\big)\le\alpha$ for every $k$, with $q^{(1)}_a=\max_kp^{(1)}_{a,k}$.

\begin{proposition}\label{prop:union}
Under (i) and (ii) of Theorem~\ref{thm:t1e}, and the conditional validity of each $p^{(2)}_k$ in (iii), the dose-level closure satisfies $\Pr(\text{claim for }a^\star\text{ and }H^\cup_{a^\star}\text{ true})\le\alpha$ for every configuration.
\end{proposition}

\begin{proof}
Let $A_0=\{a:H^\cup_a\text{ true}\}$ and fix, for each $a\in A_0$, an endpoint $k_a$ with $H_{a,k_a}$ true. On $\{a^\star\notin A_0\}$ no claim is false. As in the proof of the theorem, pass to the trial that always proceeds. On $\{a^\star=a\}$ with $a\in A_0$, the claim requires the rejection of $H^\cup_{A_0}$, hence $\cC\big(\psi_{|A_0|}(p^{(1)}_{a,k_a},(q^{(1)}_{a'})_{a'\in A_0\setminus\{a\}}),\,p^{(2)}_{k_a}\big)\le\alpha$. Since $q^{(1)}_{a'}\ge p^{(1)}_{a',k_{a'}}$ and $\psi$ and $\cC$ are non-decreasing in each argument, this implies $\cC(P_1,P_2)\le\alpha$ with $P_1=\psi_{|A_0|}\big((p^{(1)}_{a',k_{a'}})_{a'\in A_0}\big)$, which does not depend on the selected dose, is $\mathcal F_1$-measurable and is valid for the true intersection $\bigcap_{a'\in A_0}H_{a',k_{a'}}$ by Lemma~\ref{lem:comb}, and with $P_2=p^{(2)}_{k_{a^\star}}$ on $\{a^\star\in A_0\}$ and $P_2=1$ otherwise, which is conditionally valid given $\mathcal F_1$ because $H_{a^\star,k_{a^\star}}$ is true on that event. Lemma~\ref{lem:inv} completes the proof.
\end{proof}

The dose-level closure needs no Simes inequality at stage 2 and closes over $M$ doses rather than $MK$ pairs, which is why it is more powerful than the full closure (Table~\ref{tab:closure}). Its guarantee is correspondingly weaker: it bounds the probability of a false co-primary claim, not the probability of any false rejection among the $MK$ hypotheses, so it does not protect per-endpoint statements about the selected dose that are read from the same analysis. The intersection--union $p$-values $q^{(1)}_a$ (Lemma~\ref{lem:iu}) are what make its stage-1 component valid whichever dose is selected; replacing them by the same-endpoint $p$-values $p^{(1)}_{a',k_{a^\star}}$ gives the per-endpoint closure, which is valid only if the other doses are also null on $k_{a^\star}$.

\subsection{Why the per-endpoint closure fails}\label{app:counter}

Suppose each endpoint is closed separately over doses, that is, the intersection over doses $I\ni a^\star$ on endpoint $k$ is tested by $\cC\big(\psi_{|I|}((p^{(1)}_{a,k})_{a\in I}),\,p^{(2)}_k\big)$, and consider a configuration in which every dose is null on exactly one endpoint, dose $a$ on endpoint $k_a$, with $k_a\ne k_b$ for $a\ne b$ and effects on the other endpoints large enough that the corresponding $p$-values are close to zero (configurations N3 and N4 of Section~\ref{sec:sim}, with $M=3$). Then $A_0=\{1,\dots,M\}$, every claim is false, and for the selected dose $a$ and its null endpoint $k_a$ every intersection $I\ni a$ with $|I|\ge2$ contains a dose $b\ne a$ whose $p^{(1)}_{b,k_a}$ is close to zero, so that $\psi_{|I|}$ is close to zero and the intersection is rejected on endpoint $k_a$ with probability close to one. The binding condition is the elementary one, and
\begin{equation*}
\Pr(\text{false claim})\approx\sum_{a=1}^M\Pr\big(a^\star=a,\ \cC(p^{(1)}_{a,k_a},p^{(2)}_{k_a})\le\alpha\big).
\end{equation*}
If selection were independent of the stage-1 statistics of the null endpoints, each term would equal $\Pr(a^\star=a)\,\alpha$ and the sum would be $\alpha$. Selection is not independent of them: both gates favour the dose whose statistics look best in stage 1, the composite through the average of its $K$ component contrasts, of which the null endpoint is one, and the conjunctive gate through the minimum, which under this configuration is the null endpoint itself. Under the normal approximation, the stage-1 statistic $Z^{(1)}_{a,k_a}$ and the event $\{a^\star=a\}$ are positively associated, so $\Pr(a^\star=a,\,\cdot)>\Pr(a^\star=a)\Pr(\cdot)$ for each $a$ and the sum exceeds $\alpha$. For the uniform composite with $K$ equally variable, independent components the correlation between $Z^{(1)}_{a,k_a}$ and the composite contrast of dose $a$ against the other doses is $1/\sqrt{2K}$, about $0.35$ for $K=4$, diluted by the factor $v_1$ in the combination; for the conjunctive gate the ranking statistic of dose $a$ is essentially $Z^{(1)}_{a,k_a}$ itself and the association is much stronger, which is why Table~\ref{tab:closure} shows the larger inflation for that gate. Both valid closures remove the problem. In the full closure the binding intersection is $H_T$ with $T=\{(1,k_1),(2,k_2),(3,k_3)\}$, a fixed set whose stage-1 $p$-value combines three near-uniform $p$-values and is not rejected for free; in the dose-level closure the other doses enter through $q^{(1)}_b=\max_kp^{(1)}_{b,k}\ge p^{(1)}_{b,k_b}$, which is near-uniform under N3 and N4, with the same effect.

\section{Second-moment approximation for the stage-1 sample-size ratio}\label{app:bracket}

\subsection{Consistency and concordance}

Let $\hat\theta_{a,k}=\hat p_{a,k}-\hat p_{C,k}$ denote the stage-1 posterior-mean contrasts. As $n_1\to\infty$ the posterior means converge to the true probabilities and the posterior variances are of order $1/n_1$, so the composite $z$-statistic satisfies $z_a/\sqrt{n_1}\to(\bw^\top\bm\theta_a-\gamma)/\tau_a$, with $\tau_a^2$ the limit of $n_1\Var(\Delta_a\mid\bx^{(1)})$, namely the sum over the two arms of the subject-level variance of the weighted score $T$. The composite gate therefore selects, with probability tending to one, the dose maximising $(\bw^\top\bm\theta_a-\gamma)/\tau_a$. When the $\tau_a$ are equal across doses this is $a^{\mathrm C}=\arg\max_a\bw^\top\bm\theta_a$; the $\tau_a$ differ only through the dose effects on the component variances $p_{a,k}(1-p_{a,k})$, and for every configuration, correlation and weight vector in Section~\ref{sec:sim} the two maximisers coincide (the largest relative difference between the $\tau_a$ of two doses is below two per cent). Concordance, $a^{\mathrm C}=a^\circ$, is thus necessary for $\PCS\to1$, and under discordance $\PCS\to0$.

\subsection{The bracket}

Assume parallel effects $\theta_{a,k}=\theta_a$, $\delta_k=0$, uniform weights, and the normal approximation in which, for each arm $a$, $\hat p_{a,k}=p_{a,k}+\varepsilon_{a,k}$ with $\bm\varepsilon_a$ a mean-zero normal vector with common variance $\sigma^2/n_1$ and common correlation $\phi$, independent across arms. Write $\varepsilon_{a,k}=\sigma n_1^{-1/2}\big(\sqrt\phi\,U_a+\sqrt{1-\phi}\,\xi_{a,k}\big)$ with $U_a,\xi_{a,1},\dots,\xi_{a,K}$ independent standard normal, and similarly for the control with $(U_C,\xi_{C,k})$.

For the composite, $T^{\mathrm C}_a-T^{\mathrm C}_b=K^{-1}\sum_k(\hat p_{a,k}-\hat p_{b,k})$ has mean $\theta_a-\theta_b$ and variance $2\sigma^2n_1^{-1}\{\phi+(1-\phi)/K\}$, the control cancelling exactly. For the conjunctive gate, $T^{\mathrm M}_a=\min_k(\hat p_{a,k}-\hat p_{C,k})$. Ignoring the control error, $T^{\mathrm M}_a=\theta_a+\sigma n_1^{-1/2}\big(\sqrt\phi\,U_a+\sqrt{1-\phi}\min_k\xi_{a,k}\big)$, so $\E[T^{\mathrm M}_a]=\theta_a+\sigma n_1^{-1/2}\sqrt{1-\phi}\,m_K$ with $m_K=\E[\min_k\xi_{a,k}]$ common to all doses, which cancels in $\E[T^{\mathrm M}_a-T^{\mathrm M}_b]=\theta_a-\theta_b$, and $\Var[T^{\mathrm M}_a-T^{\mathrm M}_b]=2\sigma^2n_1^{-1}\{\phi+(1-\phi)V_K\}$ with $V_K=\Var[\min_k\xi_{a,k}]$. With the control error restored, $T^{\mathrm M}_a-T^{\mathrm M}_b$ acquires the term $-(e_{k_a}-e_{k_b})$, where $k_a$ and $k_b$ are the endpoints attaining the two minima and $e_k$ is the control error on endpoint $k$; this term vanishes when $k_a=k_b$ and otherwise adds variance of order $2\sigma^2n_1^{-1}(1-\phi)$, moving the factor $\phi+(1-\phi)V_K$ towards $\phi+(1-\phi)\cdot1$, the value for a single endpoint. Under the normal approximation the probability of correct selection is a function of the ratio of the pairwise signal $\theta_a-\theta_b$ to the pairwise standard deviation, so two gates reach the same $\PCS^\ast$ when $n_1$ is proportional to their variance factors; the ratio $n_1^{\mathrm C}/n_1^{\mathrm M}$ is $\{\phi+(1-\phi)/K\}/\{\phi+(1-\phi)f\}$ with $f$ between $V_K$ and $1$, which is \eqref{eq:bracket}.

The constants are $V_K=\E[\xi_{(1)}^2]-\E[\xi_{(1)}]^2$ for the minimum $\xi_{(1)}$ of $K$ independent standard normal variables, with density $K\varphi(x)\{1-\Phi(x)\}^{K-1}$: $V_2=1-1/\pi\approx0.682$, $V_3\approx0.560$ and $V_4\approx0.492$. Since $V_K$ decreases only logarithmically in $K$, $KV_K>1$ for all $K\ge2$, so the upper end of the bracket is below one whenever $\phi<1$. In the simulation study $\phi$ is evaluated at the mean control rate and a zero dose effect; the marginals there range from $0.43$ to $0.65$ across arms and endpoints, so the equal-variance assumption holds only approximately, which is one reason the bracket is wide.

\end{appendix}
